\section{Conclusion, Limitations, and Future Directions}

This paper addresses the online scheduling problem for LLM inference under memory constraints, where KV caches grow dynamically during decode and eviction risks complicate traditional scheduling approaches. We introduce the (Nested) WAIT algorithm, which provides theoretical guarantees for throughput and latency in the asymptotic regime, even when output lengths are unknown upon arrival. Our threshold-based admission control prevents eviction cascades while approaching fluid equilibrium. While our algorithms demonstrate strong performance, determining optimal threshold parameters for specific deployment scenarios remains an open problem. Furthermore, extending the algorithm to multi-GPU scenarios---prevalent in real-world applications---introduces additional complexities. Additionally, as system-level optimizations for LLM inference advance with new techniques (e.g., multi-head latent attention), developing analytical frameworks that incorporate these advances remains an important challenge. We discuss several promising directions for future work.

\subsection{More Volatile Arrivals}

In practical settings, the distribution of incoming prompts becomes non-stationary, violating the stationarity assumption underlying our model. While our algorithm can be directly adapted to scenarios with moderate variations in arrival rates—provided they remain within system capacity—by re-solving the fluid equilibrium, a thorough analysis of its performance in more volatile environments, such as during demand surges, remains an open question. Investigating how the (Nested) WAIT algorithm behaves under such conditions warrants further study.

\subsection{Extension to Multi-GPU Systems}
Our current research is limited to single-GPU systems, capable of supporting LLMs such as Llama-13B. Extending the algorithm to multi-GPU environments introduces additional considerations, including communication costs, parallelization strategies, and hardware constraints. For example, while single-GPU deployment is relatively straightforward, multi-GPU setups necessitate careful management of data parallelism, tensor parallelism, and pipeline parallelism. Moreover, communication overhead becomes a critical factor, particularly since many large-scale models require multiple GPUs for deployment. Developing a theoretical framework that incorporates these elements requires addressing communication overhead and synchronization constraints specific to distributed GPU architectures.

\subsection{Analytical Framework for Up-to-Date LLM Inference Models}
In recent years, system-level designs for LLM inference have influenced the field. For instance, the multi-head latent attention (MLA) mechanism introduced by \cite{deepseekai2024deepseekv2strongeconomicalefficient} has achieved a 90\% reduction in memory cost through KV cache compression. Integrating these state-of-the-art system-level optimizations into a theoretical framework presents a promising direction for future research.
